\clearpage\SourceLeaf{182}
\begin{center}
{\Large CORRESPONDENCE.}
\end{center}

\bigskip
\begin{center}
[TRANSCENDENTAL TRIGONOMETRY.]
\end{center}

\begin{center}\rule{2cm}{0.4pt}\end{center}

\begin{center}
[\textsc{Wachter} \emph{to} \textsc{Gauss}. \emph{Danzig, 12 December 1816.}]
\end{center}

\{.\ .\ .\ .\ .\ .\ .\ .\ Then anti-Euclidean geometry, or your geometry, would be true.
Yet its \textls{constant} remains undetermined. Why? Perhaps the
following could make that intelligible. .\ .\ .\ .\ .\ .\ .

.\ .\ .\ .\ .\ .\ .\ .\ The result of the foregoing would then be expressed as follows.

Euclidean
geometry is false. Yet true geometry must begin with
the same eleventh Euclidean axiom or with the postulate of lines and surfaces
that have the property asserted in that axiom. In place of the
straight line and plane one need only put the great circles of that sphere described with infinite radius
and its surface. There is indeed the one
inconvenience that the parts of this surface are merely symmetrical, without the
congruence they have in the plane, or that the radius on one side is
infinite and on the other imaginary. Yet it is clear how that inconvenience
is offset by the many other advantages that construction on a sphere
affords, \textls{so that perhaps even then}, if
Euclidean geometry were true, the necessity of regarding
the plane as an infinite spherical surface would disappear, but the
fruitfulness of this view might still recommend it.

Yet after I had thought all this through and had already come
quite to terms with the result, partly because I believed that the ground (la métaphysique)
of the indeterminacy necessarily inherent in geometry --- and even
